\section{The recursion theorem}
\label{sec:universal-reductions-04}

\begin{theorem}[Kleene recursion theorem]
\label{thm:recursion-theorem}
If \(f:\N\to\N\) is total computable, then there is an index \(e\) such that
\[
  \varphi_e\simeq\varphi_{f(e)}.
\]
\textit{Status: proof sketch.}
\end{theorem}

\begin{proofplan}
By universal computation and the parameter theorem, there is a total
computable function \(p\) such that
\[
  \varphi_{p(x)}(y)\simeq\varphi_{\varphi_x(x)}(y),
\]
with the right side undefined when \(\varphi_x(x)\) is undefined.  By
\cref{thm:enumeration}, let \(a\)
be an index of the total computable function \(x\mapsto f(p(x))\).  Then
\(\varphi_a(a)=f(p(a))\).  Set \(e=p(a)\).  We obtain
\[
  \varphi_e
  =\varphi_{p(a)}
  \simeq\varphi_{\varphi_a(a)}
  =\varphi_{f(p(a))}
  =\varphi_{f(e)}.
\]
This completes the fixed-point calculation using the cited universal and
parameter theorems.  The direct construction of the program transformation
\(p\) from the fixed encodings is omitted; this omission is why the recursion
theorem retains proof-sketch status.
\end{proofplan}

\begin{example}[Self-reference without a self-inspection primitive]
Suppose a total computable transformation \(f\) adds a fixed output header to
the behavior of a program whose index it receives.  The recursion theorem
supplies a program \(e\) whose behavior is exactly that of the transformed
program \(f(e)\).  The theorem does not assert \(e=f(e)\); it asserts equality
of the computed partial functions.
\end{example}

